\documentclass[../Main.tex]{subfiles}

\begin{document}
\chapter{Template details}

\begin{chapquote}{Tom Wagg (nerd)}
    \noindent Writing a thesis is so much fun!
\end{chapquote}
\chaptertoc{}

\section{General details}

Welcome to my template :) I made this for my thesis but you're very welcome to use it for yours too and edit things however you like!

The main.tex file houses the primary tex for the manuscript as you might imagine. This pulls in the styling files from the styles folder, references from the bibs folder and subchapter tex from the files folder.

You can edit individual files in the files folder and compile them separately to speed up the compilation time before you put them all together by compiling main.tex.

If you'd like to change the styling of the chapters, headers, footers or anything else it's almost all in the styles/thesis.sty folder. The styles/uw\_colours.sty file defines all of the specific colours used by UW based on \href{https://www.washington.edu/brand/brand-elements/colors/}{this website}. The styles/uw\_requirements.sty ensures that the first 3 pages match the UW requirements for theses, as listed \href{https://grad.uw.edu/current-students/enrollment-through-graduation/thesis-dissertation/}{on this page} so you probably shouldn't edit that one unless the requirements have changed.

Hope this is helpful and good luck with your thesis writing!

\begin{flushright}
    Tom Wagg (2025)
\end{flushright}

\section{Some FYIs}
The format of the thesis is set in stone for the first 3 pages, everything there you have to leave as is except for adding your details. In a crazy turn of the events (as of 2025 at least), the rest of the thesis is completely up to you. This template is my personal choice for what I think looks nice, but feel free to change things.

\textbf{Faster compiling:} Most of the time you shouldn't compile main.tex, that's going to take a \textit{while}! Instead, this is set up with subfiles, so if you compile each chapter separately, it'll still create a proper bibliography, and you can edit stuff there before you compile the whole thing.

\section{How to add a new chapter}
\begin{enumerate}
    \item Create a new file in the files folder (e.g. ``2.cogsworth.tex'') - it doesn't have to have the 2, I just do that so the chapters are in order in the files tab
    \item Copy the content from 1.introduction.tex (you want the structure to be the same)
    \item Add all of your lovely details
    \item Come back here to main.tex and add a new line with the slash subfiles command (search for 1.introduction in this file)
\end{enumerate}

\section{Your tasks}

Right - no more procrastinating - get to it!

\todo{Change your name and details on the titlepage.tex}
\todo{Add some acknowledgements and dedication}
\todo{Write your beautiful overarching introduction}
\todo{Include your awesome paper(s) as new chapters in the files folder (see above for how to)}
\todo{Find out how to add your own todos in the styles/todo\_notes.sty file (copy the slash tom command)}
\tom{Wish you luck and wave at you from the Great Postdoc Beyond}

\ifSubfilesClassLoaded{%
  \bibliographystyle{aasjournal}
  \bibliography{all}%
}{}

\end{document}
